\chapter{Geometry-Only Element Binding}

\section{Context}

After the first \texttt{Domain} transition, the library already stored
\texttt{Vertex<Dim>} as the canonical geometric entity while keeping a lazy
\texttt{Node<Dim>} cache for analysis.  That solved the ownership problem in
\texttt{Domain}, but a structural coupling remained in the element layer:
\texttt{ElementGeometry} and the concrete element families still used
\texttt{Node} as the source of physical coordinates.

Architecturally, that was still inconsistent with the ontology established in
Chapter~\ref{ch:domain-ontology}.  A \texttt{Node} is an analysis attachment
that carries degrees of freedom, PETSc numbering, and constraints.  Geometry,
however, only needs physical coordinates.  The second phase of the refactor
therefore separates these two bindings inside the element geometry layer.

\section{Problem Statement}

Before this change, a tensor-product, simplex, or serendipity element stored a
single pointer array
\[
  \{\texttt{Node<dim>*}\}_{a=1}^{n_\mathrm{node}}
\]
and used that same array both for:
\begin{enumerate}[nosep]
  \item geometric interpolation and Jacobian evaluation, and
  \item analysis-level access to nodal degrees of freedom.
\end{enumerate}

This conflated two different roles:
\begin{align*}
  \text{geometry site} &\neq \text{analysis node}.
\end{align*}

Even when both roles are collocated in space, their semantics are distinct.  A
continuum element can be fully geometrically defined from vertices without any
DOF allocation.  Conversely, an analysis formulation may decorate a point with
DOFs, constraints, or solver indices without changing its geometric role.

\section{Design Objective}

The goal of this phase was to make the geometry layer depend only on the
minimum abstraction required for isoparametric mapping:

\begin{center}
\texttt{Point<Dim>} or a type derived from it.
\end{center}

At the same time, the analysis layer still needs direct access to
\texttt{Node<Dim>} once PETSc sections and elemental DOFs are assigned.

The desired split is therefore:
\begin{align*}
  \texttt{ElementGeometry}
  &\longrightarrow
  \texttt{Point/Vertex binding} && \text{for } x(\bxi), J(\bxi), \dd\Omega, \\
  \texttt{ElementGeometry}
  &\longrightarrow
  \texttt{Node binding} && \text{for DoFs, constraints, and PETSc layout.}
\end{align*}

\section{Implemented Refactor}

\subsection{New dual binding inside concrete elements}

The three concrete element families
\begin{itemize}[nosep]
  \item \texttt{LagrangeElement},
  \item \texttt{SimplexElement}, and
  \item \texttt{SerendipityElement}
\end{itemize}
now hold two logically distinct pointer arrays:
\begin{align*}
  \texttt{points\_p} &:\ \texttt{Point<Dim>*} \\
  \texttt{nodes\_p}  &:\ \texttt{Node<Dim>*}
\end{align*}

The geometry path uses only \texttt{points\_p}.  In particular:
\begin{itemize}[nosep]
  \item \texttt{coord\_array()},
  \item \texttt{map\_local\_point()}, and
  \item \texttt{evaluate\_jacobian()}
\end{itemize}
no longer depend on \texttt{Node}.

The analysis path continues to use \texttt{nodes\_p}.  Methods such as
\texttt{node\_p()} and all element-level DOF collection code remain unchanged.

\subsection{Wrapper-level support in \texttt{ElementGeometry}}

The type-erased wrapper was extended with:
\begin{itemize}[nosep]
  \item \texttt{bind\_point(i, Point<dim>*)},
  \item \texttt{point\_p(i)},
  \item existing \texttt{bind\_node(i, Node<dim>*)}, and
  \item existing \texttt{node\_p(i)}.
\end{itemize}

This preserves the previous public API for analysis code while adding an
explicit geometry-only binding channel.

\subsection{Key compatibility rule}

Whenever a \texttt{Node} is bound, the corresponding geometry point is also
bound automatically.  In symbols:
\[
  \texttt{bind\_node}(i, n) \Longrightarrow \texttt{bind\_point}(i, n).
\]

This means:
\begin{itemize}[nosep]
  \item legacy code that only binds nodes still works unchanged;
  \item new geometry-first code can bind vertices before any DOFs exist.
\end{itemize}

\section{Domain-Level Consequences}

\subsection{Geometry-first linking}

\texttt{Domain} now exposes a dedicated
\texttt{link\_geometry\_to\_elements()} phase.  It binds each element node id to
its canonical \texttt{Vertex<dim>}:
\[
  \texttt{element.node(i)} \mapsto \texttt{vertex\_by\_id[element.node(i)]}.
\]

This enables:
\begin{itemize}[nosep]
  \item geometric mapping,
  \item Jacobian evaluation,
  \item integration point placement, and
  \item geometry-driven boundary extraction
\end{itemize}
without first materialising the analysis-node cache.

\subsection{Boundary extraction no longer needs \texttt{node\_p().id()}}

The plane-based boundary builder previously queried geometric membership
through \texttt{elem.node\_p(local).id()}, which forced the existence of node
bindings.  After the refactor, face detection uses only the canonical element
connectivity:
\[
  \texttt{elem.node(local)}.
\]

The resulting boundary element is then:
\begin{enumerate}[nosep]
  \item created from the face connectivity,
  \item immediately bound to \texttt{Vertex} objects for geometry, and
  \item optionally bound to \texttt{Node} objects if the analysis cache already
        exists.
\end{enumerate}

\subsection{Sieve assembly remains analysis-oriented}

The PETSc sieve assembly still needs analysis nodes because sieve ids and DOF
layout belong to the solver layer.  However, even there the distinction is now
cleaner:
\begin{itemize}[nosep]
  \item geometry binding prepares the element for integration and mapping;
  \item node binding prepares the element for DOF and PETSc interaction.
\end{itemize}

\section{Why this is efficient}

\subsection{No new runtime polymorphism in the hot path}

The critical loops for element integration remain unchanged:
\begin{itemize}[nosep]
  \item quadrature rule selection is still static via templates,
  \item local reference-space evaluation remains compile-time typed,
  \item Jacobian and interpolation work on fixed-size Eigen objects,
  \item the type-erased wrapper is still crossed only at the same boundary as
        before.
\end{itemize}

The only additional state is a pointer array to \texttt{Point<Dim>} inside each
concrete element.  Accessing \texttt{coord(i)} through a \texttt{Point} base
reference is not virtual; it is just a normal base-subobject access.  Therefore
this refactor does not introduce a new dynamic-dispatch cost in the Gauss loop.

\subsection{Compile-time impact}

The change slightly increases template surface area in the element headers, but
the trade-off is favourable:
\begin{itemize}[nosep]
  \item \texttt{Domain.hh} no longer needs to include Gmsh reader details;
  \item geometry-only workflows can avoid dragging analysis semantics into
        lower layers;
  \item future specialisations can target \texttt{Point}/\texttt{Vertex}
        independently of PETSc-related analysis infrastructure.
\end{itemize}

For an HPC-oriented codebase this is important because the hottest kernels
remain statically shaped, while architectural boundaries become cleaner and
easier to optimise separately.

\section{Validation}

The test suite was extended with geometry-first checks in the
\texttt{domain\_vertex\_refactor} executable:
\begin{enumerate}[nosep]
  \item a quadrilateral element can map its centre and integrate area after
        binding only vertices;
  \item a boundary element generated by
        \texttt{create\_boundary\_from\_plane()} can map and integrate from
        vertex binding alone.
\end{enumerate}

These tests validate the new invariant:
\begin{center}
\emph{element geometry is usable as soon as geometry is bound, even if analysis
nodes have not yet been materialised.}
\end{center}

\section{Architectural Interpretation}

This phase does not eliminate \texttt{Node} from the library.  It puts
\texttt{Node} back in its proper role.

\begin{itemize}[nosep]
  \item \texttt{Vertex}: canonical domain geometry/topology entity;
  \item \texttt{Point}: minimal coordinate carrier;
  \item \texttt{ElementGeometry}: geometric interpolation operator over bound
        points;
  \item \texttt{Node}: analysis decorator attached later when the model
        introduces DOFs.
\end{itemize}

That separation is especially important for future mixed and structural
formulations, where collocation between geometry sites, integration sites, and
DOF-carrying sites may occur but should not collapse the underlying concepts.

\section{Next Step}

The next logical phase is to propagate this distinction upward:
\begin{itemize}[nosep]
  \item reduce remaining read paths that only need connectivity ids but still
        access \texttt{node\_p()};
  \item make geometry-only exporters and preprocessors depend only on
        \texttt{Vertex}/\texttt{Point};
  \item move more of the solver-specific state management out of
        \texttt{Domain} and into the model/analysis side.
\end{itemize}

At that stage, \texttt{Domain} will be geometrically pure not only in storage,
but also in the majority of its operational paths.
